\documentclass[12pt]{article}
\usepackage{graphicx} % Required for inserting images


% ===================== PACKAGES =====================
\usepackage[a4paper,margin=1in]{geometry}
\usepackage{amsmath,amssymb}
\usepackage{enumitem}
\usepackage{fancyhdr}
\usepackage{xcolor} 

% ===================== HEADER & FOOTER =====================
\pagestyle{fancy}
\fancyhf{}
\lhead{CSE 400: Fundamentals of Probability in Computing}
\rhead{Lecture - 26 scribe }
\cfoot{\thepage}

% ===================== CUSTOM COMMANDS =====================
\newcommand{\solution}{
    \vspace{0.5cm}
    \noindent\textbf{\textcolor{blue}{Solution:}}
    \vspace{0.2cm}
    
    % --- STUDENT: TYPE YOUR SOLUTION BELOW THIS LINE ---
    
    \vfill 
}


% ===================== TITLE =====================
\title{
    \normalsize School of Engineering and Applied Science (SEAS), Ahmedabad University \\
    \vspace{0.2cm}
    \textbf{CSE 400: Fundamentals of Probability in Computing}\\
    \Large { Stochastic Processes: Part 1 \\ 
Strict Sense Stationary (SSS), Wide Sense Stationary (WSS) and Ergodic Process}
}
\author{} 
\date{}

\begin{document}
\maketitle

\vspace{-2cm}
\begin{center}
    \begin{tabular}{ll}
        \textbf{Name:} & {Samridhi Mehrotra \hspace{2.5in}} \\ [1.5ex]
        \textbf{Enrollment No:} & {AU2440033 \hspace{2.5in}} \\ [1.5ex]
        \textbf{Email:} & {samridhi.m@ahduni.edu.in \hspace{2.1in}} \\ [1.5ex]
        
    \end{tabular}
\end{center}

\hrule
\vspace{0.5cm}



\item \title{\textbf{}

\author{}
\date{}


\date{}


\begin{document}

\maketitle

\section*{Overview}
This lecture covers basics of stochastic processes, statistical moments, stationary processes (SSS and WSS), and ergodic processes. It includes definitions, properties, derivations, proofs, and worked examples.

\section*{Definitions and Notation}

\subsection*{Stochastic Process Definition}
A stochastic process $X(\zeta, t)$ is an entire collection of waveforms or functions over time dependent on the outcome of a random experiment.
\begin{itemize}
\item The random output heavily depends on the experiment.
\item Example: Tossing a coin gives heads or tails resulting in specific deterministic waveforms.
\end{itemize}

\subsection*{Realization}
A specific measured waveform from the experiment is a \textbf{Realization} of the process.

\subsection*{Random Variable Representation}
Fixing time $t = t_1$:
\[
X(\zeta, t_1) = X(t_1)
\]

Multiple time instants:
\[
(X(t_1), X(t_2))
\]

\subsection*{Probability Density Functions (PDF)}

\textbf{1st Order Characterization}
\[
t_1 \rightarrow f_X(x, t_1) \rightarrow F_X(x, t_1)
\]

\textbf{2nd Order Characterization}
\[
(t_1, t_2) \rightarrow f_{X_1 X_2}(x_1, x_2, t_1, t_2)
\]

\textbf{nth Order Characterization}
\[
(t_1, ..., t_n) \rightarrow f_{X_1,...,X_n}(x_1,...,x_n, t_1,...,t_n)
\]

\subsection*{Mean (First Order Moment)}
\[
\mu_X(t) = E[X(t)]
\]

\[
\mu_X(t) = \int_{-\infty}^{\infty} x \cdot f_X(x,t)\, dx
\]

\subsection*{Autocorrelation Function}
\[
R_{XX}(t_1, t_2) = E[X(t_1)X^*(t_2)]
\]

\[
R_{XX}(t_1, t_2) =
\int_{-\infty}^{\infty} \int_{-\infty}^{\infty}
x_1 x_2 f_{X_1 X_2}(x_1,x_2,t_1,t_2)\, dx_1 dx_2
\]

\subsection*{Properties of Autocorrelation}

\textbf{Conjugate Symmetry}
\[
R_{XX}(t_1,t_2) = R^*_{XX}(t_2,t_1)
\]

\textbf{Proof Sketch}
\[
R^*_{XX}(t_2,t_1) = (E[X(t_2)X^*(t_1)])^*
\]
\[
= E[X^*(t_2)X(t_1)] = R_{XX}(t_1,t_2)
\]

\textbf{Positive Definiteness}\\
Autocorrelation forms a positive definite matrix.

\subsection*{Covariance Function}
\[
C_{XX}(t_1,t_2) = E[(X(t_1)-\mu_X(t_1))(X(t_2)-\mu_X(t_2))]
\]
\[
= R_{XX}(t_1,t_2) - \mu_X(t_1)\mu_X(t_2)
\]

\subsection*{Cross-Correlation}
\[
R_{XY}(t_1,t_2) = E[X(t_1)Y(t_2)]
\]

\subsection*{Cross-Covariance}
\[
C_{XY}(t_1,t_2) = E[(X(t_1)-\mu_X(t_1))(Y(t_2)-\mu_Y(t_2))]
\]

\section*{Main Results}

\subsection*{Strict Sense Stationary (SSS) Process}

\textbf{1st Order}
\[
f_X(x,t) = f_X(x,t+\tau) = g(x)
\]
\[
\mu_X(t) = \int x g(x) dx = \mu_X = \text{constant}
\]

\textbf{2nd Order}
\[
f_X(x_1,x_2;t_1,t_2) = f_X(x_1,x_2;t_1+\tau,t_2+\tau)
\]

Choose $\tau = -t_2$:
\[
f_X(x_1,x_2;t_1,t_2) = f_X(x_1,x_2;t_1-t_2,0)
\]

\[
R_{XX}(t_1,t_2) = \phi(t_1 - t_2)
\]

\[
R_{XX}(t_1,t_2) = R_{XX}(t_1 - t_2)
\]

\subsection*{Wide Sense Stationary (WSS) Process}

\textbf{Condition 1: Constant Mean}
\[
E[X(t)] = \mu_X = \text{constant}, \quad \forall t
\]

\textbf{Condition 2: Autocorrelation depends only on lag}
\[
E[X(t_1)X^*(t_2)] = R_{XX}(\tau), \quad \tau = t_1 - t_2
\]

\subsection*{WSS vs SSS}
\[
SSS \Rightarrow WSS \quad \text{but} \quad WSS \not\Rightarrow SSS
\]

\subsection*{Gaussian Process Exception}
\[
WSS \Rightarrow SSS
\]

\subsection*{Gaussian Process Definition}
\[
f_X(x) =
\frac{1}{(2\pi)^{n/2}|R|^{1/2}}
\exp\left(-\frac{1}{2}x^T R^{-1} x\right)
\]

\[
R = E[XX^T], \quad R_{ij} = E[X(t_i)X^*(t_j)]
\]

\section*{Derivations / Proofs}

\subsection*{Positive Definite Matrix Example}

\[
A =
\begin{bmatrix}
2 & -1 & 0\\
-1 & 2 & -1\\
0 & -1 & 2
\end{bmatrix}
\]

\[
|A_1| = 2 > 0
\]

\[
|A_2| = 4 - 1 = 3 > 0
\]

\[
|A_3| = 4 > 0
\]

\section*{Worked Examples}

\subsection*{Example 1}
\[
X(t)=A\cos(\lambda t)+B\sin(\lambda t)
\]

\[
E[X(t)] = 0
\]

\[
R_{XX}(t_1,t_2) = E[A^2]\cos\lambda(t_1-t_2)
\]

\[
R_{XX}(t_1,t_2)=R_{XX}(t_1-t_2)
\]

\subsection*{Example 2}
\[
V(t)=\cos(\omega t + \theta)
\]

\[
E[V(t)] = 0
\]

\[
E[V^2(t)] = \frac{1}{2}
\]

\subsection*{HW Example: White Noise}
\[
E[X(t)] = 0
\]
\[
E[X^2(t)] = \sigma^2
\]
\[
E[X(t)X(t-\tau)] = 0 \quad (\tau \neq 0)
\]

\section*{Ergodic Process}

A stochastic process is \textbf{ergodic} if time averages equal ensemble averages.

\subsection*{Time Average}

\textbf{Mean}
\[
\langle X(t) \rangle =
\lim_{T \to \infty} \frac{1}{2T} \int_{-T}^{T} x(t)\, dt
\]

\textbf{Autocorrelation}
\[
R_{XX}(\tau)=
\lim_{T \to \infty} \frac{1}{2T}
\int_{-T}^{T} x(t)x(t+\tau)\, dt
\]

\section*{Summary}

\[
\mu_X(t)=E[X(t)]
\]
\[
R_{XX}(t_1,t_2)=E[X(t_1)X^*(t_2)]
\]
\[
C_{XX}=R_{XX}-\mu_X(t_1)\mu_X(t_2)
\]

SSS $\Rightarrow$ WSS, \quad WSS $\not\Rightarrow$ SSS

Gaussian WSS $\Rightarrow$ SSS

Ergodic: time average = ensemble average





\end{document}


